% ─── Appendix sections ───────────────────────────────────────────────────────
% \appendix is declared in main.tex before % ─── Appendix sections ───────────────────────────────────────────────────────
% \appendix is declared in main.tex before % ─── Appendix sections ───────────────────────────────────────────────────────
% \appendix is declared in main.tex before % ─── Appendix sections ───────────────────────────────────────────────────────
% \appendix is declared in main.tex before \input{sections/appendix}.

\section{Full Hyperparameter Table}
\label{app:hyperparams}

Table~\ref{tab:full_hyperparams} consolidates every configuration evaluated across
all three studies.
Parameters not listed for a config are identical to the baseline (Config~A).

\begin{table}[ht]
  \centering
  \caption{Complete hyperparameter table for all experimental configurations.
           Shared across all: $M{=}400$, $N_b{=}250$, $u_M{=}3.5$\,m/s,
           $\ell_m{=}0.75$\,m, $\gamma_M{=}\pi/6$, $T_s{=}0.02$\,s,
           $N_{\mathrm{opt}}{=}1500$, $\alpha_{\mathrm{launch}}{=}35^\circ$,
           ball ($m{=}0.0577$\,kg, $r{=}0.0327$\,m).
           ``Strat'' = stratified exploration; ``Rnd'' = random.}
  \label{tab:full_hyperparams}
  \small
  \begin{tabular}{llcccccccl}
    \toprule
    Config & Sub-project & $z_{\mathrm{rel}}$ & $\ell_M$ & $T$ & $\ell_c$ &
    $\ell_s^{\mathrm{init}}$ & $N_{\mathrm{exp}}$ & Exploration & Result \\
    \midrule
    A               & mc-pilot/           & 0.5\,m & 1.75\,m & 0.70\,s & 0.5\,m & [1.0, 1.0] & 5  & Rnd & 5/5 hits \\
    B               & mc-pilot-elevated/  & 1.0\,m & 1.90\,m & 0.85\,s & 0.5\,m & [1.0, 1.0] & 5  & Rnd & 0/10 (seed=1); 10/10 (seed=2) \\
    B-Strat         & mc-pilot-elevated/  & 1.0\,m & 1.90\,m & 0.85\,s & 0.5\,m & [1.0, 1.0] & 5  & Strat $[0, u_M]$ & 10/10 \\
    C-Strat         & mc-pilot-elevated/  & 1.5\,m & 2.15\,m & 0.95\,s & 0.5\,m & [1.0, 1.0] & 5  & Strat $[0, u_M]$ & 10/10 \\
    D-Strat         & mc-pilot-elevated/  & 2.0\,m & 2.35\,m & 1.00\,s & 0.5\,m & [1.0, 1.0] & 5  & Strat $[0, u_M]$ & 10/10 \\
    E-Strat         & mc-pilot-elevated/  & 0.0\,m & 1.10\,m & 0.55\,s & 0.5\,m & [1.0, 1.0] & 5  & Strat $[0, u_M]$ & 0/10 \\
    E-Strat2        & mc-pilot-elevated/  & 0.0\,m & 1.10\,m & 0.55\,s & 0.5\,m & [1.0, 1.0] & 5  & Strat $[2.5, 3.5]$ & 0/10 \\
    E-Strat3        & mc-pilot-elevated/  & 0.0\,m & 1.10\,m & 0.55\,s & 0.5\,m & [1.0, 1.0] & 10 & Strat $[2.5, 3.5]$ & 0/10 \\
    E-Strat4        & mc-pilot-elevated/  & 0.0\,m & 1.10\,m & 0.55\,s & 0.25\,m & [1.0, 1.0] & 10 & Strat $[2.5, 3.5]$ & 0/10 \\
    E-Strat5        & mc-pilot-elevated/  & 0.0\,m & 1.10\,m & 0.55\,s & 0.25\,m & [0.15, 0.15] & 10 & Strat $[2.5, 3.5]$ & 10/10 \\
    PBE-B           & mc-pilot-pb-elevated/ & 1.0\,m & 1.90\,m & 0.85\,s & 0.5\,m & [1.0, 1.0] & 5 & Strat $[0, u_M]$ & \needdata{run pending} \\
    PBE-C           & mc-pilot-pb-elevated/ & 1.5\,m & 2.15\,m & 0.95\,s & 0.5\,m & [1.0, 1.0] & 5 & Strat $[0, u_M]$ & \needdata{run pending} \\
    PBE-D           & mc-pilot-pb-elevated/ & 2.0\,m & 2.35\,m & 1.00\,s & 0.5\,m & [1.0, 1.0] & 5 & Strat $[0, u_M]$ & \needdata{run pending} \\
    \bottomrule
  \end{tabular}
\end{table}

\section{Zero-Mean Unlearnability: Proof Sketch}
\label{app:zeromean}

\begin{proposition}[Zero-mean unlearnability, restated]
Let $\boldsymbol{\epsilon} \sim \mathcal{N}(\mathbf{0}, \Sigma_\epsilon)$ be
additive noise on the commanded release velocity.
The optimal policy under noise-awareness is identical to the noiseless optimal policy.
\end{proposition}

\begin{proof}[Proof sketch]
Let $f : \mathbb{R}^3 \to \mathbb{R}^2$ denote the deterministic landing map
(position as a function of release velocity), assumed smooth and invertible near the
optimum.
By first-order Taylor expansion around $\mathbf{v}^*$:
\[
  f(\mathbf{v}^* + \boldsymbol{\epsilon}) \approx f(\mathbf{v}^*) + J_f \boldsymbol{\epsilon},
  \qquad J_f = \tfrac{\partial f}{\partial \mathbf{v}}\big|_{\mathbf{v}^*}.
\]
The expected landing position is
\[
  \mathbb{E}_\epsilon[f(\mathbf{v}^* + \boldsymbol{\epsilon})]
  = f(\mathbf{v}^*) + J_f \mathbb{E}[\boldsymbol{\epsilon}]
  = f(\mathbf{v}^*),
\]
since $\mathbb{E}[\boldsymbol{\epsilon}] = \mathbf{0}$.
The noise-averaged cost gradient with respect to the policy output $\mathbf{v}^*$ is
\[
  \nabla_{\mathbf{v}^*} J_{\mathrm{aware}}
  = \nabla_{\mathbf{v}^*} \mathbb{E}_\epsilon[c(f(\mathbf{v}^* + \boldsymbol{\epsilon}))]
  = J_f^\top \nabla_{\mathbf{p}} c\bigl(f(\mathbf{v}^*)\bigr),
\]
which is identical to $\nabla_{\mathbf{v}^*} J_{\mathrm{naive}}$ computed with
$\boldsymbol{\epsilon} = \mathbf{0}$.
Therefore, the gradient is zero at the same $\mathbf{v}^*$ for both the aware and
naive objectives, and the optimal policies coincide.
The noise increases the variance of $c$ (and thus the honest cost value) but does
not shift the optimal aim.
\end{proof}

\begin{remark}
The result relies on linearity of $f$ near the optimum and zero mean of
$\boldsymbol{\epsilon}$.
For skewed or heavy-tailed noise distributions the approximation breaks and a
genuine gap between aware and naive policies can emerge.
\end{remark}

\section{Lengthscale Scaling Rule Derivation}
\label{app:lengthscale}

The RBF policy (Eq.~\ref{eq:policy}) uses a squared-exponential kernel with
lengthscale $\ell_s$.
The sensitivity between two target positions $\mathbf{P}_1$ and $\mathbf{P}_2$ at
Euclidean distance $d$ is
\[
  S(d,\, \ell_s) = \exp\!\Bigl(-\tfrac{d^2}{2\ell_s^2}\Bigr).
\]
Config A (target range $\approx 1.0$\,m, $\ell_s = 1.0$\,m) achieves
sensitivity $S_A = \exp(-0.5) \approx 0.607$ between the two extremes of its domain.

For a new configuration with target range $R$, we require the same sensitivity:
\[
  S(R,\, \ell_s^{\mathrm{new}}) \stackrel{!}{=} S_A = 0.607.
\]
Setting $\exp(-R^2 / (2\ell_s^2)) = \exp(-0.5)$ gives
\[
  \frac{R^2}{2\ell_s^2} = \frac{1}{2}
  \implies \ell_s = R.
\]
Hence $\ell_s = R$ for all target ranges---but in practice, the paper's fixed
$\ell_s = 1.0$\,m is used even when $R \ll 1.0$\,m.
The issue is that Config A has $R \approx 1.0$\,m, so $\ell_s = R$ happens to hold
by coincidence.
For Config E with $R = 0.35$\,m, setting $\ell_s = 0.35$\,m recovers
$S = 0.607$.
In our experiments, $\ell_s = 0.15$\,m (which gives $S = 0.066$, lower than 0.607)
was found to work well because even sharper sensitivity promotes faster learning when
the GP has limited exploration data.
The conservative recommendation is $\ell_s = R$; the empirical finding is that
$\ell_s \approx 0.15 R$ is safe and often faster.

\section{Particle Freeze Approximation at Landing}
\label{app:freeze}

The GP propagates $M$ particles forward in time.
Once a particle's $z$-coordinate satisfies $z \leq 0$, its state is frozen at the
last step above ground and the cost is evaluated at that frozen landing position for
all remaining timesteps.

This introduces a systematic over-estimation of the horizontal landing distance
proportional to the residual horizontal velocity at the frozen step.
For timestep $T_s = 0.02$\,s and a typical vertical velocity of 2\,m/s at impact,
the true ground crossing occurs $\Delta t = z_{\mathrm{last}} / v_z \approx 0.01$\,s
before the frozen step, during which the ball travels an additional
$\Delta x \approx v_x \cdot \Delta t \approx 2\text{--}4$\,cm horizontally
(for $v_x \approx 2\text{--}4$\,m/s).

In actual rollouts (NumPy and PyBullet), landing is detected at $z \leq 0$ and
refined by linear interpolation to the exact $z = 0$ crossing, so training data are
not affected by this error.
The approximation applies only to the Monte Carlo particles used for the gradient
estimate inside \texttt{apply\_policy}, where implementing per-particle
interpolation would break batched tensor operations.
The resulting cost estimates are slightly pessimistic (the policy is penalised for
landing a few centimetres beyond the target when the true landing is closer), but
the bias is consistent across particles and policy updates, so optimisation still
converges to the correct policy.
Reducing $T_s$ or implementing a soft landing indicator (a smooth sigmoid of $z$
rather than a hard step) would eliminate this error at the cost of increased
computation.
.

\section{Full Hyperparameter Table}
\label{app:hyperparams}

Table~\ref{tab:full_hyperparams} consolidates every configuration evaluated across
all three studies.
Parameters not listed for a config are identical to the baseline (Config~A).

\begin{table}[ht]
  \centering
  \caption{Complete hyperparameter table for all experimental configurations.
           Shared across all: $M{=}400$, $N_b{=}250$, $u_M{=}3.5$\,m/s,
           $\ell_m{=}0.75$\,m, $\gamma_M{=}\pi/6$, $T_s{=}0.02$\,s,
           $N_{\mathrm{opt}}{=}1500$, $\alpha_{\mathrm{launch}}{=}35^\circ$,
           ball ($m{=}0.0577$\,kg, $r{=}0.0327$\,m).
           ``Strat'' = stratified exploration; ``Rnd'' = random.}
  \label{tab:full_hyperparams}
  \small
  \begin{tabular}{llcccccccl}
    \toprule
    Config & Sub-project & $z_{\mathrm{rel}}$ & $\ell_M$ & $T$ & $\ell_c$ &
    $\ell_s^{\mathrm{init}}$ & $N_{\mathrm{exp}}$ & Exploration & Result \\
    \midrule
    A               & mc-pilot/           & 0.5\,m & 1.75\,m & 0.70\,s & 0.5\,m & [1.0, 1.0] & 5  & Rnd & 5/5 hits \\
    B               & mc-pilot-elevated/  & 1.0\,m & 1.90\,m & 0.85\,s & 0.5\,m & [1.0, 1.0] & 5  & Rnd & 0/10 (seed=1); 10/10 (seed=2) \\
    B-Strat         & mc-pilot-elevated/  & 1.0\,m & 1.90\,m & 0.85\,s & 0.5\,m & [1.0, 1.0] & 5  & Strat $[0, u_M]$ & 10/10 \\
    C-Strat         & mc-pilot-elevated/  & 1.5\,m & 2.15\,m & 0.95\,s & 0.5\,m & [1.0, 1.0] & 5  & Strat $[0, u_M]$ & 10/10 \\
    D-Strat         & mc-pilot-elevated/  & 2.0\,m & 2.35\,m & 1.00\,s & 0.5\,m & [1.0, 1.0] & 5  & Strat $[0, u_M]$ & 10/10 \\
    E-Strat         & mc-pilot-elevated/  & 0.0\,m & 1.10\,m & 0.55\,s & 0.5\,m & [1.0, 1.0] & 5  & Strat $[0, u_M]$ & 0/10 \\
    E-Strat2        & mc-pilot-elevated/  & 0.0\,m & 1.10\,m & 0.55\,s & 0.5\,m & [1.0, 1.0] & 5  & Strat $[2.5, 3.5]$ & 0/10 \\
    E-Strat3        & mc-pilot-elevated/  & 0.0\,m & 1.10\,m & 0.55\,s & 0.5\,m & [1.0, 1.0] & 10 & Strat $[2.5, 3.5]$ & 0/10 \\
    E-Strat4        & mc-pilot-elevated/  & 0.0\,m & 1.10\,m & 0.55\,s & 0.25\,m & [1.0, 1.0] & 10 & Strat $[2.5, 3.5]$ & 0/10 \\
    E-Strat5        & mc-pilot-elevated/  & 0.0\,m & 1.10\,m & 0.55\,s & 0.25\,m & [0.15, 0.15] & 10 & Strat $[2.5, 3.5]$ & 10/10 \\
    PBE-B           & mc-pilot-pb-elevated/ & 1.0\,m & 1.90\,m & 0.85\,s & 0.5\,m & [1.0, 1.0] & 5 & Strat $[0, u_M]$ & \needdata{run pending} \\
    PBE-C           & mc-pilot-pb-elevated/ & 1.5\,m & 2.15\,m & 0.95\,s & 0.5\,m & [1.0, 1.0] & 5 & Strat $[0, u_M]$ & \needdata{run pending} \\
    PBE-D           & mc-pilot-pb-elevated/ & 2.0\,m & 2.35\,m & 1.00\,s & 0.5\,m & [1.0, 1.0] & 5 & Strat $[0, u_M]$ & \needdata{run pending} \\
    \bottomrule
  \end{tabular}
\end{table}

\section{Zero-Mean Unlearnability: Proof Sketch}
\label{app:zeromean}

\begin{proposition}[Zero-mean unlearnability, restated]
Let $\boldsymbol{\epsilon} \sim \mathcal{N}(\mathbf{0}, \Sigma_\epsilon)$ be
additive noise on the commanded release velocity.
The optimal policy under noise-awareness is identical to the noiseless optimal policy.
\end{proposition}

\begin{proof}[Proof sketch]
Let $f : \mathbb{R}^3 \to \mathbb{R}^2$ denote the deterministic landing map
(position as a function of release velocity), assumed smooth and invertible near the
optimum.
By first-order Taylor expansion around $\mathbf{v}^*$:
\[
  f(\mathbf{v}^* + \boldsymbol{\epsilon}) \approx f(\mathbf{v}^*) + J_f \boldsymbol{\epsilon},
  \qquad J_f = \tfrac{\partial f}{\partial \mathbf{v}}\big|_{\mathbf{v}^*}.
\]
The expected landing position is
\[
  \mathbb{E}_\epsilon[f(\mathbf{v}^* + \boldsymbol{\epsilon})]
  = f(\mathbf{v}^*) + J_f \mathbb{E}[\boldsymbol{\epsilon}]
  = f(\mathbf{v}^*),
\]
since $\mathbb{E}[\boldsymbol{\epsilon}] = \mathbf{0}$.
The noise-averaged cost gradient with respect to the policy output $\mathbf{v}^*$ is
\[
  \nabla_{\mathbf{v}^*} J_{\mathrm{aware}}
  = \nabla_{\mathbf{v}^*} \mathbb{E}_\epsilon[c(f(\mathbf{v}^* + \boldsymbol{\epsilon}))]
  = J_f^\top \nabla_{\mathbf{p}} c\bigl(f(\mathbf{v}^*)\bigr),
\]
which is identical to $\nabla_{\mathbf{v}^*} J_{\mathrm{naive}}$ computed with
$\boldsymbol{\epsilon} = \mathbf{0}$.
Therefore, the gradient is zero at the same $\mathbf{v}^*$ for both the aware and
naive objectives, and the optimal policies coincide.
The noise increases the variance of $c$ (and thus the honest cost value) but does
not shift the optimal aim.
\end{proof}

\begin{remark}
The result relies on linearity of $f$ near the optimum and zero mean of
$\boldsymbol{\epsilon}$.
For skewed or heavy-tailed noise distributions the approximation breaks and a
genuine gap between aware and naive policies can emerge.
\end{remark}

\section{Lengthscale Scaling Rule Derivation}
\label{app:lengthscale}

The RBF policy (Eq.~\ref{eq:policy}) uses a squared-exponential kernel with
lengthscale $\ell_s$.
The sensitivity between two target positions $\mathbf{P}_1$ and $\mathbf{P}_2$ at
Euclidean distance $d$ is
\[
  S(d,\, \ell_s) = \exp\!\Bigl(-\tfrac{d^2}{2\ell_s^2}\Bigr).
\]
Config A (target range $\approx 1.0$\,m, $\ell_s = 1.0$\,m) achieves
sensitivity $S_A = \exp(-0.5) \approx 0.607$ between the two extremes of its domain.

For a new configuration with target range $R$, we require the same sensitivity:
\[
  S(R,\, \ell_s^{\mathrm{new}}) \stackrel{!}{=} S_A = 0.607.
\]
Setting $\exp(-R^2 / (2\ell_s^2)) = \exp(-0.5)$ gives
\[
  \frac{R^2}{2\ell_s^2} = \frac{1}{2}
  \implies \ell_s = R.
\]
Hence $\ell_s = R$ for all target ranges---but in practice, the paper's fixed
$\ell_s = 1.0$\,m is used even when $R \ll 1.0$\,m.
The issue is that Config A has $R \approx 1.0$\,m, so $\ell_s = R$ happens to hold
by coincidence.
For Config E with $R = 0.35$\,m, setting $\ell_s = 0.35$\,m recovers
$S = 0.607$.
In our experiments, $\ell_s = 0.15$\,m (which gives $S = 0.066$, lower than 0.607)
was found to work well because even sharper sensitivity promotes faster learning when
the GP has limited exploration data.
The conservative recommendation is $\ell_s = R$; the empirical finding is that
$\ell_s \approx 0.15 R$ is safe and often faster.

\section{Particle Freeze Approximation at Landing}
\label{app:freeze}

The GP propagates $M$ particles forward in time.
Once a particle's $z$-coordinate satisfies $z \leq 0$, its state is frozen at the
last step above ground and the cost is evaluated at that frozen landing position for
all remaining timesteps.

This introduces a systematic over-estimation of the horizontal landing distance
proportional to the residual horizontal velocity at the frozen step.
For timestep $T_s = 0.02$\,s and a typical vertical velocity of 2\,m/s at impact,
the true ground crossing occurs $\Delta t = z_{\mathrm{last}} / v_z \approx 0.01$\,s
before the frozen step, during which the ball travels an additional
$\Delta x \approx v_x \cdot \Delta t \approx 2\text{--}4$\,cm horizontally
(for $v_x \approx 2\text{--}4$\,m/s).

In actual rollouts (NumPy and PyBullet), landing is detected at $z \leq 0$ and
refined by linear interpolation to the exact $z = 0$ crossing, so training data are
not affected by this error.
The approximation applies only to the Monte Carlo particles used for the gradient
estimate inside \texttt{apply\_policy}, where implementing per-particle
interpolation would break batched tensor operations.
The resulting cost estimates are slightly pessimistic (the policy is penalised for
landing a few centimetres beyond the target when the true landing is closer), but
the bias is consistent across particles and policy updates, so optimisation still
converges to the correct policy.
Reducing $T_s$ or implementing a soft landing indicator (a smooth sigmoid of $z$
rather than a hard step) would eliminate this error at the cost of increased
computation.
.

\section{Full Hyperparameter Table}
\label{app:hyperparams}

Table~\ref{tab:full_hyperparams} consolidates every configuration evaluated across
all three studies.
Parameters not listed for a config are identical to the baseline (Config~A).

\begin{table}[ht]
  \centering
  \caption{Complete hyperparameter table for all experimental configurations.
           Shared across all: $M{=}400$, $N_b{=}250$, $u_M{=}3.5$\,m/s,
           $\ell_m{=}0.75$\,m, $\gamma_M{=}\pi/6$, $T_s{=}0.02$\,s,
           $N_{\mathrm{opt}}{=}1500$, $\alpha_{\mathrm{launch}}{=}35^\circ$,
           ball ($m{=}0.0577$\,kg, $r{=}0.0327$\,m).
           ``Strat'' = stratified exploration; ``Rnd'' = random.}
  \label{tab:full_hyperparams}
  \small
  \begin{tabular}{llcccccccl}
    \toprule
    Config & Sub-project & $z_{\mathrm{rel}}$ & $\ell_M$ & $T$ & $\ell_c$ &
    $\ell_s^{\mathrm{init}}$ & $N_{\mathrm{exp}}$ & Exploration & Result \\
    \midrule
    A               & mc-pilot/           & 0.5\,m & 1.75\,m & 0.70\,s & 0.5\,m & [1.0, 1.0] & 5  & Rnd & 5/5 hits \\
    B               & mc-pilot-elevated/  & 1.0\,m & 1.90\,m & 0.85\,s & 0.5\,m & [1.0, 1.0] & 5  & Rnd & 0/10 (seed=1); 10/10 (seed=2) \\
    B-Strat         & mc-pilot-elevated/  & 1.0\,m & 1.90\,m & 0.85\,s & 0.5\,m & [1.0, 1.0] & 5  & Strat $[0, u_M]$ & 10/10 \\
    C-Strat         & mc-pilot-elevated/  & 1.5\,m & 2.15\,m & 0.95\,s & 0.5\,m & [1.0, 1.0] & 5  & Strat $[0, u_M]$ & 10/10 \\
    D-Strat         & mc-pilot-elevated/  & 2.0\,m & 2.35\,m & 1.00\,s & 0.5\,m & [1.0, 1.0] & 5  & Strat $[0, u_M]$ & 10/10 \\
    E-Strat         & mc-pilot-elevated/  & 0.0\,m & 1.10\,m & 0.55\,s & 0.5\,m & [1.0, 1.0] & 5  & Strat $[0, u_M]$ & 0/10 \\
    E-Strat2        & mc-pilot-elevated/  & 0.0\,m & 1.10\,m & 0.55\,s & 0.5\,m & [1.0, 1.0] & 5  & Strat $[2.5, 3.5]$ & 0/10 \\
    E-Strat3        & mc-pilot-elevated/  & 0.0\,m & 1.10\,m & 0.55\,s & 0.5\,m & [1.0, 1.0] & 10 & Strat $[2.5, 3.5]$ & 0/10 \\
    E-Strat4        & mc-pilot-elevated/  & 0.0\,m & 1.10\,m & 0.55\,s & 0.25\,m & [1.0, 1.0] & 10 & Strat $[2.5, 3.5]$ & 0/10 \\
    E-Strat5        & mc-pilot-elevated/  & 0.0\,m & 1.10\,m & 0.55\,s & 0.25\,m & [0.15, 0.15] & 10 & Strat $[2.5, 3.5]$ & 10/10 \\
    PBE-B           & mc-pilot-pb-elevated/ & 1.0\,m & 1.90\,m & 0.85\,s & 0.5\,m & [1.0, 1.0] & 5 & Strat $[0, u_M]$ & \needdata{run pending} \\
    PBE-C           & mc-pilot-pb-elevated/ & 1.5\,m & 2.15\,m & 0.95\,s & 0.5\,m & [1.0, 1.0] & 5 & Strat $[0, u_M]$ & \needdata{run pending} \\
    PBE-D           & mc-pilot-pb-elevated/ & 2.0\,m & 2.35\,m & 1.00\,s & 0.5\,m & [1.0, 1.0] & 5 & Strat $[0, u_M]$ & \needdata{run pending} \\
    \bottomrule
  \end{tabular}
\end{table}

\section{Zero-Mean Unlearnability: Proof Sketch}
\label{app:zeromean}

\begin{proposition}[Zero-mean unlearnability, restated]
Let $\boldsymbol{\epsilon} \sim \mathcal{N}(\mathbf{0}, \Sigma_\epsilon)$ be
additive noise on the commanded release velocity.
The optimal policy under noise-awareness is identical to the noiseless optimal policy.
\end{proposition}

\begin{proof}[Proof sketch]
Let $f : \mathbb{R}^3 \to \mathbb{R}^2$ denote the deterministic landing map
(position as a function of release velocity), assumed smooth and invertible near the
optimum.
By first-order Taylor expansion around $\mathbf{v}^*$:
\[
  f(\mathbf{v}^* + \boldsymbol{\epsilon}) \approx f(\mathbf{v}^*) + J_f \boldsymbol{\epsilon},
  \qquad J_f = \tfrac{\partial f}{\partial \mathbf{v}}\big|_{\mathbf{v}^*}.
\]
The expected landing position is
\[
  \mathbb{E}_\epsilon[f(\mathbf{v}^* + \boldsymbol{\epsilon})]
  = f(\mathbf{v}^*) + J_f \mathbb{E}[\boldsymbol{\epsilon}]
  = f(\mathbf{v}^*),
\]
since $\mathbb{E}[\boldsymbol{\epsilon}] = \mathbf{0}$.
The noise-averaged cost gradient with respect to the policy output $\mathbf{v}^*$ is
\[
  \nabla_{\mathbf{v}^*} J_{\mathrm{aware}}
  = \nabla_{\mathbf{v}^*} \mathbb{E}_\epsilon[c(f(\mathbf{v}^* + \boldsymbol{\epsilon}))]
  = J_f^\top \nabla_{\mathbf{p}} c\bigl(f(\mathbf{v}^*)\bigr),
\]
which is identical to $\nabla_{\mathbf{v}^*} J_{\mathrm{naive}}$ computed with
$\boldsymbol{\epsilon} = \mathbf{0}$.
Therefore, the gradient is zero at the same $\mathbf{v}^*$ for both the aware and
naive objectives, and the optimal policies coincide.
The noise increases the variance of $c$ (and thus the honest cost value) but does
not shift the optimal aim.
\end{proof}

\begin{remark}
The result relies on linearity of $f$ near the optimum and zero mean of
$\boldsymbol{\epsilon}$.
For skewed or heavy-tailed noise distributions the approximation breaks and a
genuine gap between aware and naive policies can emerge.
\end{remark}

\section{Lengthscale Scaling Rule Derivation}
\label{app:lengthscale}

The RBF policy (Eq.~\ref{eq:policy}) uses a squared-exponential kernel with
lengthscale $\ell_s$.
The sensitivity between two target positions $\mathbf{P}_1$ and $\mathbf{P}_2$ at
Euclidean distance $d$ is
\[
  S(d,\, \ell_s) = \exp\!\Bigl(-\tfrac{d^2}{2\ell_s^2}\Bigr).
\]
Config A (target range $\approx 1.0$\,m, $\ell_s = 1.0$\,m) achieves
sensitivity $S_A = \exp(-0.5) \approx 0.607$ between the two extremes of its domain.

For a new configuration with target range $R$, we require the same sensitivity:
\[
  S(R,\, \ell_s^{\mathrm{new}}) \stackrel{!}{=} S_A = 0.607.
\]
Setting $\exp(-R^2 / (2\ell_s^2)) = \exp(-0.5)$ gives
\[
  \frac{R^2}{2\ell_s^2} = \frac{1}{2}
  \implies \ell_s = R.
\]
Hence $\ell_s = R$ for all target ranges---but in practice, the paper's fixed
$\ell_s = 1.0$\,m is used even when $R \ll 1.0$\,m.
The issue is that Config A has $R \approx 1.0$\,m, so $\ell_s = R$ happens to hold
by coincidence.
For Config E with $R = 0.35$\,m, setting $\ell_s = 0.35$\,m recovers
$S = 0.607$.
In our experiments, $\ell_s = 0.15$\,m (which gives $S = 0.066$, lower than 0.607)
was found to work well because even sharper sensitivity promotes faster learning when
the GP has limited exploration data.
The conservative recommendation is $\ell_s = R$; the empirical finding is that
$\ell_s \approx 0.15 R$ is safe and often faster.

\section{Particle Freeze Approximation at Landing}
\label{app:freeze}

The GP propagates $M$ particles forward in time.
Once a particle's $z$-coordinate satisfies $z \leq 0$, its state is frozen at the
last step above ground and the cost is evaluated at that frozen landing position for
all remaining timesteps.

This introduces a systematic over-estimation of the horizontal landing distance
proportional to the residual horizontal velocity at the frozen step.
For timestep $T_s = 0.02$\,s and a typical vertical velocity of 2\,m/s at impact,
the true ground crossing occurs $\Delta t = z_{\mathrm{last}} / v_z \approx 0.01$\,s
before the frozen step, during which the ball travels an additional
$\Delta x \approx v_x \cdot \Delta t \approx 2\text{--}4$\,cm horizontally
(for $v_x \approx 2\text{--}4$\,m/s).

In actual rollouts (NumPy and PyBullet), landing is detected at $z \leq 0$ and
refined by linear interpolation to the exact $z = 0$ crossing, so training data are
not affected by this error.
The approximation applies only to the Monte Carlo particles used for the gradient
estimate inside \texttt{apply\_policy}, where implementing per-particle
interpolation would break batched tensor operations.
The resulting cost estimates are slightly pessimistic (the policy is penalised for
landing a few centimetres beyond the target when the true landing is closer), but
the bias is consistent across particles and policy updates, so optimisation still
converges to the correct policy.
Reducing $T_s$ or implementing a soft landing indicator (a smooth sigmoid of $z$
rather than a hard step) would eliminate this error at the cost of increased
computation.
.

\section{Full Hyperparameter Table}
\label{app:hyperparams}

Table~\ref{tab:full_hyperparams} consolidates every configuration evaluated across
all three studies.
Parameters not listed for a config are identical to the baseline (Config~A).

\begin{table}[ht]
  \centering
  \caption{Complete hyperparameter table for all experimental configurations.
           Shared across all: $M{=}400$, $N_b{=}250$, $u_M{=}3.5$\,m/s,
           $\ell_m{=}0.75$\,m, $\gamma_M{=}\pi/6$, $T_s{=}0.02$\,s,
           $N_{\mathrm{opt}}{=}1500$, $\alpha_{\mathrm{launch}}{=}35^\circ$,
           ball ($m{=}0.0577$\,kg, $r{=}0.0327$\,m).
           ``Strat'' = stratified exploration; ``Rnd'' = random.}
  \label{tab:full_hyperparams}
  \small
  \begin{tabular}{llcccccccl}
    \toprule
    Config & Sub-project & $z_{\mathrm{rel}}$ & $\ell_M$ & $T$ & $\ell_c$ &
    $\ell_s^{\mathrm{init}}$ & $N_{\mathrm{exp}}$ & Exploration & Result \\
    \midrule
    A               & mc-pilot/           & 0.5\,m & 1.75\,m & 0.70\,s & 0.5\,m & [1.0, 1.0] & 5  & Rnd & 5/5 hits \\
    B               & mc-pilot-elevated/  & 1.0\,m & 1.90\,m & 0.85\,s & 0.5\,m & [1.0, 1.0] & 5  & Rnd & 0/10 (seed=1); 10/10 (seed=2) \\
    B-Strat         & mc-pilot-elevated/  & 1.0\,m & 1.90\,m & 0.85\,s & 0.5\,m & [1.0, 1.0] & 5  & Strat $[0, u_M]$ & 10/10 \\
    C-Strat         & mc-pilot-elevated/  & 1.5\,m & 2.15\,m & 0.95\,s & 0.5\,m & [1.0, 1.0] & 5  & Strat $[0, u_M]$ & 10/10 \\
    D-Strat         & mc-pilot-elevated/  & 2.0\,m & 2.35\,m & 1.00\,s & 0.5\,m & [1.0, 1.0] & 5  & Strat $[0, u_M]$ & 10/10 \\
    E-Strat         & mc-pilot-elevated/  & 0.0\,m & 1.10\,m & 0.55\,s & 0.5\,m & [1.0, 1.0] & 5  & Strat $[0, u_M]$ & 0/10 \\
    E-Strat2        & mc-pilot-elevated/  & 0.0\,m & 1.10\,m & 0.55\,s & 0.5\,m & [1.0, 1.0] & 5  & Strat $[2.5, 3.5]$ & 0/10 \\
    E-Strat3        & mc-pilot-elevated/  & 0.0\,m & 1.10\,m & 0.55\,s & 0.5\,m & [1.0, 1.0] & 10 & Strat $[2.5, 3.5]$ & 0/10 \\
    E-Strat4        & mc-pilot-elevated/  & 0.0\,m & 1.10\,m & 0.55\,s & 0.25\,m & [1.0, 1.0] & 10 & Strat $[2.5, 3.5]$ & 0/10 \\
    E-Strat5        & mc-pilot-elevated/  & 0.0\,m & 1.10\,m & 0.55\,s & 0.25\,m & [0.15, 0.15] & 10 & Strat $[2.5, 3.5]$ & 10/10 \\
    PBE-B           & mc-pilot-pb-elevated/ & 1.0\,m & 1.90\,m & 0.85\,s & 0.5\,m & [1.0, 1.0] & 5 & Strat $[0, u_M]$ & \needdata{run pending} \\
    PBE-C           & mc-pilot-pb-elevated/ & 1.5\,m & 2.15\,m & 0.95\,s & 0.5\,m & [1.0, 1.0] & 5 & Strat $[0, u_M]$ & \needdata{run pending} \\
    PBE-D           & mc-pilot-pb-elevated/ & 2.0\,m & 2.35\,m & 1.00\,s & 0.5\,m & [1.0, 1.0] & 5 & Strat $[0, u_M]$ & \needdata{run pending} \\
    \bottomrule
  \end{tabular}
\end{table}

\section{Zero-Mean Unlearnability: Proof Sketch}
\label{app:zeromean}

\begin{proposition}[Zero-mean unlearnability, restated]
Let $\boldsymbol{\epsilon} \sim \mathcal{N}(\mathbf{0}, \Sigma_\epsilon)$ be
additive noise on the commanded release velocity.
The optimal policy under noise-awareness is identical to the noiseless optimal policy.
\end{proposition}

\begin{proof}[Proof sketch]
Let $f : \mathbb{R}^3 \to \mathbb{R}^2$ denote the deterministic landing map
(position as a function of release velocity), assumed smooth and invertible near the
optimum.
By first-order Taylor expansion around $\mathbf{v}^*$:
\[
  f(\mathbf{v}^* + \boldsymbol{\epsilon}) \approx f(\mathbf{v}^*) + J_f \boldsymbol{\epsilon},
  \qquad J_f = \tfrac{\partial f}{\partial \mathbf{v}}\big|_{\mathbf{v}^*}.
\]
The expected landing position is
\[
  \mathbb{E}_\epsilon[f(\mathbf{v}^* + \boldsymbol{\epsilon})]
  = f(\mathbf{v}^*) + J_f \mathbb{E}[\boldsymbol{\epsilon}]
  = f(\mathbf{v}^*),
\]
since $\mathbb{E}[\boldsymbol{\epsilon}] = \mathbf{0}$.
The noise-averaged cost gradient with respect to the policy output $\mathbf{v}^*$ is
\[
  \nabla_{\mathbf{v}^*} J_{\mathrm{aware}}
  = \nabla_{\mathbf{v}^*} \mathbb{E}_\epsilon[c(f(\mathbf{v}^* + \boldsymbol{\epsilon}))]
  = J_f^\top \nabla_{\mathbf{p}} c\bigl(f(\mathbf{v}^*)\bigr),
\]
which is identical to $\nabla_{\mathbf{v}^*} J_{\mathrm{naive}}$ computed with
$\boldsymbol{\epsilon} = \mathbf{0}$.
Therefore, the gradient is zero at the same $\mathbf{v}^*$ for both the aware and
naive objectives, and the optimal policies coincide.
The noise increases the variance of $c$ (and thus the honest cost value) but does
not shift the optimal aim.
\end{proof}

\begin{remark}
The result relies on linearity of $f$ near the optimum and zero mean of
$\boldsymbol{\epsilon}$.
For skewed or heavy-tailed noise distributions the approximation breaks and a
genuine gap between aware and naive policies can emerge.
\end{remark}

\section{Lengthscale Scaling Rule Derivation}
\label{app:lengthscale}

The RBF policy (Eq.~\ref{eq:policy}) uses a squared-exponential kernel with
lengthscale $\ell_s$.
The sensitivity between two target positions $\mathbf{P}_1$ and $\mathbf{P}_2$ at
Euclidean distance $d$ is
\[
  S(d,\, \ell_s) = \exp\!\Bigl(-\tfrac{d^2}{2\ell_s^2}\Bigr).
\]
Config A (target range $\approx 1.0$\,m, $\ell_s = 1.0$\,m) achieves
sensitivity $S_A = \exp(-0.5) \approx 0.607$ between the two extremes of its domain.

For a new configuration with target range $R$, we require the same sensitivity:
\[
  S(R,\, \ell_s^{\mathrm{new}}) \stackrel{!}{=} S_A = 0.607.
\]
Setting $\exp(-R^2 / (2\ell_s^2)) = \exp(-0.5)$ gives
\[
  \frac{R^2}{2\ell_s^2} = \frac{1}{2}
  \implies \ell_s = R.
\]
Hence $\ell_s = R$ for all target ranges---but in practice, the paper's fixed
$\ell_s = 1.0$\,m is used even when $R \ll 1.0$\,m.
The issue is that Config A has $R \approx 1.0$\,m, so $\ell_s = R$ happens to hold
by coincidence.
For Config E with $R = 0.35$\,m, setting $\ell_s = 0.35$\,m recovers
$S = 0.607$.
In our experiments, $\ell_s = 0.15$\,m (which gives $S = 0.066$, lower than 0.607)
was found to work well because even sharper sensitivity promotes faster learning when
the GP has limited exploration data.
The conservative recommendation is $\ell_s = R$; the empirical finding is that
$\ell_s \approx 0.15 R$ is safe and often faster.

\section{Particle Freeze Approximation at Landing}
\label{app:freeze}

The GP propagates $M$ particles forward in time.
Once a particle's $z$-coordinate satisfies $z \leq 0$, its state is frozen at the
last step above ground and the cost is evaluated at that frozen landing position for
all remaining timesteps.

This introduces a systematic over-estimation of the horizontal landing distance
proportional to the residual horizontal velocity at the frozen step.
For timestep $T_s = 0.02$\,s and a typical vertical velocity of 2\,m/s at impact,
the true ground crossing occurs $\Delta t = z_{\mathrm{last}} / v_z \approx 0.01$\,s
before the frozen step, during which the ball travels an additional
$\Delta x \approx v_x \cdot \Delta t \approx 2\text{--}4$\,cm horizontally
(for $v_x \approx 2\text{--}4$\,m/s).

In actual rollouts (NumPy and PyBullet), landing is detected at $z \leq 0$ and
refined by linear interpolation to the exact $z = 0$ crossing, so training data are
not affected by this error.
The approximation applies only to the Monte Carlo particles used for the gradient
estimate inside \texttt{apply\_policy}, where implementing per-particle
interpolation would break batched tensor operations.
The resulting cost estimates are slightly pessimistic (the policy is penalised for
landing a few centimetres beyond the target when the true landing is closer), but
the bias is consistent across particles and policy updates, so optimisation still
converges to the correct policy.
Reducing $T_s$ or implementing a soft landing indicator (a smooth sigmoid of $z$
rather than a hard step) would eliminate this error at the cost of increased
computation.
